\chapter*{Abstract}

\iffalse{
The abstract should be a concise summary of what you have done.
It should cover, in roughly 1 to 2 sentences per topic:
\begin{itemize}
	\item The domain, i.e., explain the application domain in which your thesis is applicable.
	\item The problem, i.e., motivate what the problem is you are solving.
	\item The method, i.e., describe the method you have proposed, and, compare it, briefly \& high level to other approaches. In particular, focus on the benefits of your approach versus the other approach(es)
	\item The evaluation, i.e., how did you evaluate your method, what results did you obtain  and what do the results indicate?
\end{itemize}

}\fi



% The domain 
Life Cycle Assessment is an environmental reporting framework in which the input and output flows of a process are quantified in the Life Cycle Inventory analysis (LCI analysis) phase.
Waste flows are of particular interest for sustainability.

% The problem
LCI analysis has the potential to be automated using object-centric process mining, but existing process mining algorithms focus on the control flow of a process, in which events represent synchronization points in the life cycle of objects.
An LCI analysis, however, requires knowledge about the material flow, which interprets events as transformations of input objects to output objects.


% The method 
To discover a model of the material flow from an object-centric event log, which we identify with the object flow, we propose the Object Flow Miner, a discovery algorithm that constructs Object Flow Nets (OFNs).
These visualize the object flow through quantified transformations of input objects into output objects for each activity.
Knowledge about the inputs and outputs of events can be used to compute waste flows, and to detect implicit waste by calculating the difference between input and output mass.


% The evaluation 
The algorithm was evaluated on public object-centric event logs, where it discovered accurate as-is models of the object flow or derived suitable process models through filtering.
We demonstrated our approach on the hinge production log, where it semi-automated the LCI analysis by providing a quantified OFN containing information about the waste flows.
